% Pregunta 340 (2020-enero-OEP2019, n.º 7): un prestamista sigue siéndolo cuando sube r.
% delta=0,8; dotación (6,2); r=0,25 -> 0,6; E0=(4,22;4,22), E1=(4,03;5,16). Preferencia revelada.
% Autor: Víctor Gutiérrez Marcos
\documentclass[tikz,border=6pt]{standalone}
\input{tikz-tcee}
\begin{document}
\begin{tikzpicture}[tcee,x=0.62cm,y=0.45cm]
  \ejes{$C_1$}{$C_2$}{9.4}{12.8}
  % Conjunto presupuestario antiguo con C1 > w1 (planes de prestatario)
  \fill[area] (6,0) -- (6,2) -- (7.6,0) -- cycle;
  % Curvas de indiferencia
  \draw[curva3] plot[domain=2.15:9,samples=60] ({\x},{exp((2.5927-ln(\x))/0.8)});
  \draw[curva3,desplazada] plot[domain=2.25:9,samples=60] ({\x},{exp((2.7054-ln(\x))/0.8)});
  % Rectas presupuestarias
  \draw[curva1] (0,9.5) -- (7.6,0) node[etiqueta,below right=0pt and -2pt] {$r$};
  \draw[curvanueva] (0,11.6) -- (7.25,0) node[etiqueta,below left=0pt and -2pt] {$r'$};
  % Dotación
  \draw[guia] (6,0) node[below] {$\omega_1$} -- (6,2) -- (0,2) node[left] {$\omega_2$};
  \node[punto,label={[etiqueta]below left:$\omega$}] at (6,2) {};
  % Óptimos
  \draw[guia] (4.22,0) node[below] {$C_1^*$} -- (4.22,4.22);
  \node[punto,fill=tceemorado,label={[etiqueta,text=tceemorado]below left:$E_0$}] at (4.22,4.22) {};
  \node[punto,fill=tceerojo,label={[etiqueta,text=tceerojo]above right:$E_1$}] at (4.03,5.16) {};
  % Nota
  \node[etiqueta,font=\footnotesize,align=left,anchor=west] (n) at (5.5,7.2) {planes de prestatario\\con $r'$: ya eran\\asequibles con $r$\\y los rechazó};
  \draw[->,thin,tceegris] (n.south) ++(-0.2,0.1) -- (6.62,1.0);
\end{tikzpicture}
\end{document}
